\section{依存语法详解}

\subsection{依存语法的基本概念}

\begin{figure}[H]
\centering
\includegraphics[width=0.95\textwidth]{slides-images/slide-014.jpg}
\caption{依存结构示例：``Bills on ports and immigration were submitted by Senator Brownback, Republican of Kansas''\protect\footnotemark}
\end{figure}
\footnotetext{Slides 第14页。}

依存语法的核心元素：

\begin{importantbox}{依存语法的形式化定义}
\begin{itemize}
    \item \textbf{依存关系}（Dependency）：一个有向弧，从\textbf{中心词}（head/governor）指向\textbf{依存词}（dependent/modifier）
    \item \textbf{依存树}（Dependency Tree）：所有依存关系构成一棵有向树，有唯一的根节点（root）
    \item \textbf{关系标签}（Relation Labels）：每条弧上标注语法关系类型，如 nsubj（名词主语）、obj（宾语）、amod（形容词修饰语）、det（限定词）、obl（斜格）等
    \item \textbf{虚拟根节点}（ROOT）：通常在句子前添加一个虚拟的 ROOT 节点，其唯一依存词是句子的主要动词
\end{itemize}
\end{importantbox}

\begin{figure}[H]
\centering
\includegraphics[width=0.95\textwidth]{slides-images/slide-015.jpg}
\caption{带有关系标签的依存结构：每条弧上标注了具体的语法关系类型\protect\footnotemark}
\end{figure}
\footnotetext{Slides 第15页。}

\subsection{依存语法的悠久历史}

\begin{figure}[H]
\centering
\includegraphics[width=0.95\textwidth]{slides-images/slide-018.jpg}
\caption{依存语法的历史：从 Panini 到现代\protect\footnotemark}
\end{figure}
\footnotetext{Slides 第18页。}

\begin{knowledgebox}{依存语法的历史脉络}
依存语法是人类历史上最古老的句法形式化方法：
\begin{itemize}
    \item \textbf{Panini}（约公元前 4-8 世纪）：最早的语法学家，用依存关系描述梵语（Sanskrit）的语法结构。他的语法最初以\textbf{口述方式}传承数百年，令人惊叹
    \item \textbf{阿拉伯语法传统}（第一千年）：使用依存语法描述阿拉伯语
    \item \textbf{Reed-Kellogg 句子图解}：在美国教育中广泛使用的句子图解法，本质上是依存语法的一种变体
    \item \textbf{Lucien Tesni\`ere}（20世纪）：现代依存语法的奠基人，确立了从中心词指向依存词的箭头方向
    \item 相比之下，\textbf{短语结构语法/CFG} 是非常晚近的发明（1940s--1950s，由 Chomsky 等人系统化）
\end{itemize}
\end{knowledgebox}

Manning 特别提到了一个有趣的事实：计算机科学家通常从 Chomsky 层次结构了解 Chomsky，但 Chomsky 层次结构最初并非为了形式语言理论——而是为了论证人类语言不能用正则/有限状态文法来充分描述。

\subsection{箭头方向的约定}

在依存语法中，箭头方向存在两种约定：

\begin{itemize}
    \item \textbf{Head $\rightarrow$ Dependent}：Tesni\`ere 传统，也是本课程使用的方向
    \item \textbf{Dependent $\rightarrow$ Head}：另一些研究者使用的方向
\end{itemize}

\begin{warningbox}{注意箭头方向约定}
在阅读不同的依存语法文献时，必须首先确认所使用的箭头方向约定。不同的工具和标注方案可能采用不同的约定，混淆方向会导致对结构的完全误读。
\end{warningbox}

\subsection{依存分析的决策依据}

Manning 总结了依存分析器在做出依存决策时需要考虑的四类信息：

\begin{enumerate}
    \item \textbf{词汇亲和性（Bilexical Affinities）}：两个词之间是否存在合理的依存关系。例如 ``discussion $\leftarrow$ issues'' 是合理的，但 ``the $\leftarrow$ completed'' 则不合理
    \item \textbf{依存距离（Dependency Distance）}：大多数依存关系发生在\textbf{相邻或近距离}的词之间。长距离依存虽然存在，但较为少见
    \item \textbf{中间成分（Intervening Material）}：依存弧很少跨越动词或标点符号
    \item \textbf{中心词配价（Valency of Heads）}：中心词能接受多少个依存词。例如 ``broke'' 通常有一个主语（左侧）和一个可选的宾语（右侧），但不能接受多个宾语
\end{enumerate}

\begin{figure}[H]
\centering
\includegraphics[width=0.95\textwidth]{slides-images/slide-020.jpg}
\caption{依存分析的决策信息来源\protect\footnotemark}
\end{figure}
\footnotetext{Slides 第20页。}

\subsection{投射性与非投射性（Projectivity）}

\begin{figure}[H]
\centering
\includegraphics[width=0.95\textwidth]{slides-images/slide-021.jpg}
\caption{非投射性依存示例：依存弧交叉\protect\footnotemark}
\end{figure}
\footnotetext{Slides 第21页。}

\begin{importantbox}{投射性与非投射性依存}
\begin{itemize}
    \item \textbf{投射性（Projective）}：依存弧不交叉，等价于上下文无关文法可生成的结构。大多数依存关系满足投射性
    \item \textbf{非投射性（Non-projective）}：依存弧存在交叉。例如 ``I'll give a talk tomorrow on neural networks'' 中，``on neural networks'' 修饰 ``talk''，``tomorrow'' 修饰 ``give''，两条弧交叉
\end{itemize}
疑问句中的 Wh-移位也常产生非投射依存：``Who did Bill buy the coffee from yesterday?'' 中 ``who'' 是 ``from'' 的宾语，但被移到了句首。
\end{importantbox}

本课程的作业中只处理投射性依存分析。

\subsection{本章小结}

依存语法通过中心词-依存词关系描述句子结构，具有悠久的历史传统。依存树具有唯一根节点和树形结构约束。分析器的决策依赖词汇亲和性、距离偏好、中间成分和配价约束。大多数依存关系是投射性的，但自然语言中也存在非投射性的情况。

%% ========== 树库 ==========

